\documentclass[11pt]{article}

\usepackage[margin=1.25in]{geometry}
\usepackage{setspace}
\usepackage{microtype}
\usepackage{amsmath}
\usepackage{amssymb}
\usepackage{hyperref}
\usepackage{natbib}
\usepackage{titlesec}

\setstretch{1.25}
\setlength{\parindent}{0pt}
\setlength{\parskip}{0.8em}
\title{Identity After Authenticity:\\
Physiological Realism, Invariant Structure, and the Social Representation}

\author{Flyxion}

\date{\today}

\begin{document}

\maketitle
\begin{abstract}

This essay begins with a moral injunction encountered in adolescence: \textit{Esse quam videri}—to be rather than to seem. What initially appears as a character ideal is reinterpreted as a structural problem that anticipates a central tension in modern social theory. In \textit{The Society of the Spectacle}, Guy Debord diagnoses a historical condition in which being is displaced by appearing and social relations are mediated by images. His critique unfolds across semiotics, political economy, and ontology, culminating in a narrative of lost authenticity.

The present study reframes that diagnosis by grounding the distinction between being and seeming in the physiology of perception and the mathematics of invariance. The visual system does not passively record images but learns to detect transformation-invariant structure—curvature, depth, shadow geometry, and material reflectance. Sparse coding renders cognition metabolically efficient yet structurally vulnerable to camouflage and mimicry. Under this interpretation, the spectacle functions as a socio-cognitive analogue of biological camouflage: it exploits compression constraints to simulate invariants without preserving generative substrate.

By introducing a divergence functional that measures the gap between social signals and their structural or energetic basis, the argument shifts from moral exhortation to formal constraint. The imperative to be rather than to seem becomes a demand for structural congruence between representation and generative process. The essay concludes by situating this framework alongside the theory of profilicity, which rejects nostalgia for an unmediated past while recognizing that identity is technologically scaffolded. The aim is not a recovery of pre-representational immediacy, but the preservation of coherence under transformation.

\end{abstract}

\newpage
\section{Introduction: The Tension of Being and Seeming}

When I began high school, I noticed that its motto was the Latin phrase \textit{Esse quam videri}. The phrase translates as ``to be rather than to seem.'' At first glance, it reads as a simple moral directive. It encourages sincerity over hypocrisy, substance over display. Yet the structure of the phrase reveals a deeper philosophical tension. It posits a distinction between ontological presence and perceptual appearance. It implies that seeming can detach from being, that appearance can simulate what it does not structurally possess.

Long before encountering critical social theory, this motto established a conceptual axis: the difference between structural reality and performed representation. That axis would later find a systematic articulation in the work of \textit{The Society of the Spectacle}, where the displacement of being by appearing is described not as a personal failing but as a civilizational condition.

The tension between being and seeming is not merely ethical. It is epistemological and ontological. It concerns the conditions under which perception tracks structure. To seem is to generate signals within a perceptual field. To be is to maintain structural coherence under transformation. The distinction is therefore not reducible to sincerity or intention. It concerns invariance.

In modernity, the proliferation of mediated representation intensifies this tension. Images circulate independently of the processes they depict. Reputation precedes action. Branding precedes production. Evaluation precedes experience. The perceptual order is saturated with signals optimized for recognition rather than correspondence. The question raised by \textit{Esse quam videri} is whether signal and substrate remain congruent.

This question becomes explicit in the diagnostic framework offered by Guy Debord in 1967. His account does not begin with physiology or mathematics, but with political economy. He argues that twentieth-century capitalism has evolved into a regime in which images become the primary commodity. Social life is reorganized around display. The spectacle is not a collection of media artifacts but a social relation mediated by images.

The spectacle, in this sense, represents a structural inversion. Where prior social formations organized value around material production, the spectacle organizes value around visibility. Being is subordinated to appearing. Economic and social participation require continuous presentation. The show becomes the substance. The image acquires autonomy.

Yet Debord frames this inversion as a historical loss. He suggests that a prior unity of life has dissolved into representation. According to this narrative, what was once directly lived has become mere display. The moral undertone echoes the motto \textit{Esse quam videri}. Authentic presence has been displaced by staged visibility.

The present essay does not simply repeat that critique. Instead, it reframes the tension through a physiological and mathematical lens. The central claim is that the distinction between being and seeming can be modeled in terms of structural invariance and signal distortion. Perception evolved to detect transformation-invariant features in the environment. When signals simulate invariance without preserving substrate structure, cognitive systems are misled. The spectacle can therefore be interpreted as a large-scale exploitation of perceptual compression.

The shift proposed here moves from socio-economic critique to structural modeling. Rather than asking whether authenticity has been lost, we ask whether invariant structure is preserved under transformation. The difference is subtle but decisive. The demand is not nostalgic return but measurable congruence.

In what follows, the theory of the spectacle will be reconstructed in its original Marxist and ontological context. It will then be reinterpreted through the physiology of perception, sparse coding theory, and transformation groups. The moral injunction to be rather than to seem will become a formal constraint on divergence between signal and substrate. The question will not be whether appearance exists, but whether appearance tracks structural reality.

The tension introduced by a school motto thus becomes the organizing principle of a broader inquiry. It is an inquiry into how perception learns physics, how compression enables vulnerability, and how social systems can amplify that vulnerability into a regime of spectacle.

\section{Ciceronian Origins: Civic Virtue and the Priority of Being}

The phrase \textit{Esse quam videri} predates modern social theory by nearly two millennia. It appears in Cicero’s \textit{De Amicitia} (\textit{On Friendship}), where it functions as a civic and moral principle rather than a perceptual or structural one. For Cicero, the injunction to be rather than to seem is a warning against the cultivation of reputation detached from character. The Roman statesman frames friendship and public life as domains in which integrity must precede display. Virtue is not theatrical performance but stable disposition.

In its classical context, the phrase does not reject appearance altogether. Roman political culture depended heavily on visibility, rhetoric, and honor. Rather, Cicero insists that visible honor must correspond to genuine moral substance. Reputation without virtue is fragile; it collapses under scrutiny. The distinction between being and seeming therefore functions as a stabilizing constraint on public life. Appearance must be grounded in ethical structure.

The phrase later acquired institutional embodiment. It was adopted in 1893 as the official motto of the state of North Carolina. In this civic usage, \textit{Esse quam videri} becomes a declaration of collective identity. It signals a commitment to integrity at the level of governance and public character. The motto’s endurance suggests that the tension between substance and performance is not a modern invention but a persistent feature of social organization.

Within the present framework, the Ciceronian origin acquires renewed significance. Cicero’s concern is moral congruence. Debord’s concern is socio-economic displacement of lived reality by representation. The physiological account developed earlier concerns invariant structure under transformation. These three domains converge on a shared principle: appearance must track substrate.

In Cicero’s formulation, the substrate is ethical character. In Debord’s analysis, the substrate is lived experience and material production. In the physiological model, the substrate is transformation-invariant structure in physical and energetic systems. The divergence functional $D(S,B)$ generalizes Cicero’s moral concern into a structural one. If signal $S$—whether reputation, brand, or profile—deviates from substrate $B$—whether character, labor, or energetic flow—then seeming dominates being.

The Ciceronian maxim thus anticipates both the critique of spectacle and its reformulation. It does not deny the necessity of appearance. Public life requires visibility. What it demands is congruence. Virtue must be invariant under scrutiny. It must remain stable across transformations of context, illumination, and audience. In modern terms, it requires that the compressed representation $\phi(S)$ correspond to the invariant structure of $B$.

The historical endurance of \textit{Esse quam videri} indicates that the tension between substance and performance is not reducible to digital media or capitalist spectacle. What changes across epochs is the scale and velocity at which divergence can be amplified. In a regime of high-throughput signaling, where profiles and brands circulate globally, the cost of structural incongruence increases. Compression intensifies. Camouflage becomes systemic.

To return to the phrase, then, is not to indulge in nostalgia. It is to recognize that the demand for congruence between being and seeming has long served as a stabilizing principle of civic life. The present analysis extends that principle beyond ethics into physiology and formal modeling. It treats Cicero’s injunction not as mere moral advice, but as an early articulation of a constraint on signal–substrate divergence.

In this sense, \textit{Esse quam videri} becomes a bridge across domains. It links Roman civic virtue, Marxist social critique, perceptual neuroscience, and mathematical invariance. Its persistence underscores the structural depth of the problem. To be rather than to seem is not simply to act sincerely. It is to preserve coherence under transformation in systems increasingly optimized for display.

\section{Biblical Interiorization: 1 Samuel 16:7 and the Primacy of the Heart}

The distinction between being and seeming also appears in the Hebrew Bible. In 1 Samuel 16:7, the prophet Samuel is instructed during the anointing of David: ``For the Lord sees not as man sees: man looks on the outward appearance, but the Lord looks on the heart.'' The verse relocates the tension from civic virtue to theological anthropology. It asserts that external form is an unreliable indicator of inner structure. Divine perception penetrates surface display.

The biblical formulation sharpens the contrast. Human perception is described as appearance-bound. It privileges visible markers: stature, strength, social presence. The divine gaze, by contrast, accesses interior disposition. The heart in Hebrew thought does not denote mere emotion but the seat of intention, will, and moral structure. Being is defined not by outward visibility but by invariant interior alignment.

Placed alongside Cicero, the verse introduces an epistemic asymmetry. Cicero assumes that moral congruence can be recognized within civic life through consistent conduct. 1 Samuel suggests that surface evaluation is structurally limited. Human vision remains vulnerable to misclassification. The theological claim anticipates the physiological argument advanced earlier: perception operates under compression constraints. It extracts sparse features and is therefore susceptible to camouflage.

In the present framework, the biblical distinction may be reformulated without reducing its theological force. Let outward appearance correspond to signal $S$, and let the heart correspond to structural substrate $B$. Human observers rely on compressed representations $\phi(S)$ to classify character. Divine perception, in the biblical narrative, bypasses this compression, accessing $B$ directly. The divergence functional
\[
D(S,B) = \|\phi(S) - \phi(\sigma(B))\|
\]
is invisible to ordinary observers when $\phi(S)$ successfully simulates invariant cues. The theological claim is that ultimate judgment does not rely on $\phi(S)$ but on $B$ itself.

This contrast reinforces the earlier argument. Sparse coding implies that classification systems depend on minimal basis features. In evolutionary contexts, this suffices for survival. In moral and social contexts, however, it opens the possibility of systemic distortion. The spectacle can amplify signals optimized to match recognition templates while concealing substrate incoherence. The biblical warning acknowledges the structural gap between appearance and interior structure.

Importantly, 1 Samuel 16:7 does not eliminate appearance. David himself must still be seen, anointed, and recognized. Rather, the verse reorders epistemic priority. Outward markers do not guarantee inward congruence. The heart functions as invariant substrate under transformation. Integrity, in this sense, is structural stability across shifting contexts.

The convergence of Cicero and 1 Samuel reveals a deep continuity across cultural traditions. Both insist that surface display is insufficient as a measure of reality. Both imply that congruence between signal and substrate is the condition of legitimacy. The modern spectacle intensifies the stakes by multiplying channels of display and accelerating feedback loops. Yet the underlying tension remains ancient.

In physiological terms, the heart of the matter is invariance. Whether framed theologically, civically, or mathematically, the principle is consistent: being must persist under transformation. When outward appearance is decoupled from interior structure, seeming overtakes being. When interior structure remains stable across contexts, appearance becomes faithful projection rather than camouflage.

Thus, the injunction of 1 Samuel 16:7 complements \textit{Esse quam videri}. Both articulate the primacy of structural congruence. Both warn against mistaking optimized signal for grounded substance. Within a society saturated by spectacle and profilicity, their shared insight becomes not merely devotional or moral, but structurally diagnostic.

\section{General Principle and Generating Function: Beyond Particular Implementations}

The injunction \textit{Esse quam videri}, when read alongside Cicero and 1 Samuel 16:7, can be extended beyond moral psychology into a structural epistemology. To be rather than to seem does not merely contrast inner sincerity with outer display. It can also be interpreted as a methodological demand: to seek the generating principle rather than fixate on its specific implementation.

In mathematics, a generating function encodes an entire class of objects through a compact structural rule. Particular instances may vary, yet they are derivable from a shared functional form. The generating principle remains invariant across manifestations. The specific implementation is a projection of that deeper structure under particular boundary conditions.

Similarly, in physics, symmetries and conservation laws operate as general constraints from which observable phenomena are derived. The curvature of spacetime generates gravitational trajectories; the Lagrangian generates equations of motion. One does not mistake the specific trajectory for the underlying variational principle. The latter persists across transformations of coordinates and context.

Applied to perception, this distinction parallels the difference between invariant detection and surface feature recognition. The organism does not merely memorize specific pixel arrays. It learns transformation-invariant structure. It identifies curvature, continuity, and lawful change. The generating structure of the environment becomes the basis for reliable inference. Surface variation is secondary.

The spectacle, by contrast, often promotes attachment to specific implementations. It elevates the instantiated image over the generating structure. A brand becomes detached from the production process that sustains it. A profile becomes detached from the energetic and relational substrate that enables it. The visible implementation substitutes for the underlying generative dynamics.

To read \textit{Esse quam videri} in structural terms is therefore to privilege the generating function over its projection. Being corresponds to the invariant rule or structural constraint that gives rise to appearances. Seeming corresponds to a particular instantiation optimized for recognition within a perceptual or social field. When the generating structure is coherent, particular appearances vary without compromising integrity. When the generating structure is hollow, implementations must be continually refined to preserve recognition.

Formally, let $G$ denote a generating operator acting on a parameter space $\Theta$, producing observable instances $x = G(\theta)$. The spectacle privileges $x$ while neglecting $G$. Structural congruence requires that the observable instance faithfully reflect the constraints encoded in $G$. Divergence arises when $x$ is engineered to trigger recognition independently of the underlying generative rule.

In cognitive terms, this corresponds to overfitting surface features while neglecting causal structure. In social terms, it corresponds to optimizing signals while eroding substrate. In moral terms, it corresponds to cultivating reputation while neglecting character. Across domains, the pattern is consistent: attention shifts from the rule to the projection.

Thus, to be rather than to seem is to orient inquiry toward invariant generating structure. It is to ask what transformation law gives rise to observed patterns. It is to test whether specific implementations remain coherent under variation. Invariance under transformation becomes the criterion of authenticity.

This reframing also protects against nostalgia. One need not posit a lost pure implementation. Instead, one evaluates whether the generating function itself remains stable. If identity is technologically mediated, then the question becomes whether the identity-generating processes preserve structural reciprocity. If economic systems are mediated by images, the question becomes whether image production reflects underlying energetic constraints.

The injunction therefore becomes methodological as well as ethical. Seek the rule, not merely the display. Evaluate invariance, not merely instance. In a regime where spectacle and profilicity amplify particular implementations, the discipline of returning to generating structure restores congruence. Being is not a static essence. It is the persistence of a generative law across transformation.

\section{Guy Debord and the Theory of the Spectacle}

In 1967, Guy Debord published \textit{The Society of the Spectacle}, a work that sought to diagnose the structural logic of advanced capitalism. Writing within a Marxist framework and as a member of the Situationist International, Debord argued that the dominant mode of production had undergone a qualitative shift. The production of material goods no longer defined the core of economic life. Instead, the commodification and circulation of images became central. Social relations were increasingly mediated not by direct interaction, but by representation.

Debord’s most cited formulation defines the spectacle as ``a social relationship between people that is mediated by images.'' The spectacle is not reducible to television, advertising, or entertainment, though these are its most visible manifestations. It is a totalizing condition in which appearing displaces being as the organizing principle of value. The image becomes autonomous. What matters is not the structural substrate of production, but the perceptual field in which commodities, institutions, and identities are staged.

The Marxist foundation of Debord’s argument remains explicit. Classical industrial capitalism, as analyzed by Marx, alienated workers from the products of their labor. In spectacular capitalism, alienation extends further. Individuals are separated from lived experience itself. Movement becomes tourism, sexuality becomes pornography, information becomes infotainment. In each case, activity is transformed into display. The commodity no longer merely satisfies needs; it contemplates itself in a world of its own making.

Debord’s theory rests on three interlocking pillars: semiotics, political economy, and ontology.

At the semiotic level, the spectacle perfects the separation between sign and referent. Influenced by structuralist and post-structuralist thinkers, Debord argues that representations detach from the material realities they once indexed. The sign no longer points back to a stable thing; it circulates within a system of differential relations. Brands, celebrities, and media events function as self-referential sign clusters. Meaning emerges within discourse rather than through correspondence. The copy becomes more powerful than the original.

At the level of political economy, the spectacle operates as a mode of production. It is self-generating and self-perpetuating. Value is determined by visibility. Economic flows follow attention gradients. Hierarchies emerge from control over representation rather than control over physical resources alone. Debord describes the spectacle as capital accumulated to the point where it becomes image. In this formulation, capital no longer merely organizes labor; it organizes perception.

Ontologically, the spectacle functions as what Debord calls an ``appearance machine.'' It blurs the distinction between reality and illusion. The shift he describes is not merely from having to appearing, but from being to appearing. Appearance acquires normative authority. That which appears is validated; that which remains invisible is structurally marginalized. The spectacle reconstructs religious illusion in material form. It produces a secularized transcendence in which images assume the role once occupied by sacred symbols.

From these three pillars emerges Debord’s narrative of lost authenticity. He famously writes that all that was once directly lived has become mere representation. The former unity of life is declared irretrievable. Leisure becomes performance. Stardom transforms individuals into images. Communication becomes pseudo-dialogue. The spectacle generates a condition in which human beings encounter each other primarily as mediated constructs.

Alienation thus becomes systemic. The spectacle alienates individuals not only from labor but from movement, sexuality, information, and history. It constructs a world of visibility in which participation requires constant display. The individual is driven into forms of compensatory fantasy. Consumption substitutes for dialogue. Recognition substitutes for reciprocity.

Debord’s response to this diagnosis is revolutionary. He calls for the construction of situations in which individuals collectively reappropriate their lived experience. He envisions a form of social organization in which representation no longer dominates action. Yet his account remains tied to a narrative of historical loss. The spectacle is treated as a deviation from a more authentic mode of being.

The present essay accepts the structural insight of Debord’s diagnosis while suspending its nostalgia. Rather than presuming an original unity of life, we reinterpret the spectacle through the physiology of perception and the mathematics of invariance. The question becomes not whether authenticity once existed, but how perceptual systems detect structure and how social systems exploit those detection mechanisms.

In this reframing, the spectacle is not only a socio-economic phenomenon but a cognitive one. It is a regime that operates by simulating invariants recognized by compressed perceptual systems. To understand how appearing can displace being, we must first understand how organisms learn to distinguish invariant structure from transient surface variation. The next section turns to that physiological foundation.

\section{Physiological Realism and the Physics of Perception}

The distinction between being and seeming is often treated as a moral or metaphysical contrast. Yet it is also grounded in the physiology of perception. The visual system does not passively record images. It performs structured inference. The retina and downstream cortical circuits extract features that correspond to stable regularities in the external world. In doing so, the organism learns to detect invariants.

Light arriving at the retina is subject to transformation. Objects rotate, translate, scale, and vary under changing illumination. Shadows shift as light sources move. Surfaces occlude one another. Reflections distort environmental structure according to material properties. Despite these variations, the organism must identify stable entities. Perception therefore privileges features that remain constant across transformation groups.

Depth perception illustrates this principle. Binocular disparity arises because each eye receives a slightly different projection of the same scene. The brain integrates these differences to infer three-dimensional structure. Curvature is not directly visible as a property in itself; it is inferred from gradients of shading, contour deformation, and motion parallax. Shadows encode information about spatial topology, yet their position shifts with illumination. The organism learns to treat certain transformations as irrelevant while preserving others as structurally significant.

To formalize this, let $X$ represent the space of environmental states and let $G$ denote a transformation group acting on $X$. Transformations in $G$ may include rotations, translations, scaling operations, and illumination shifts. A perceptual feature $f : X \to \mathbb{R}^n$ is invariant under $G$ if
\[
f(g \cdot x) = f(x) \quad \text{for all } g \in G, \; x \in X.
\]
Such invariants correspond to stable physical constraints. They define what it means for something to be, in the perceptual sense. Being is that which maintains structural coherence under transformation.

However, the organism cannot encode the full environmental state. Metabolic and computational constraints necessitate compression. Neural systems employ sparse coding strategies, representing stimuli through a limited set of basis features. Let $\phi(x)$ denote a compressed representation of $x \in X$ in a lower-dimensional space $Z$. Then
\[
\phi(x) = \sum_{i=1}^{k} a_i b_i,
\]
where the basis $\{b_i\}$ is selected to minimize representational cost while preserving task-relevant structure. Sparsity ensures efficiency, but it also creates vulnerability.

Because classification relies on a reduced feature set, it can be manipulated. Camouflage in nature exploits precisely this constraint. If an organism can produce surface features such that its compressed representation approximates that of its environment, detection fails. Formally, if for prey state $x$ and background state $y$,
\[
\|\phi(x) - \phi(y)\| < \epsilon,
\]
for a discrimination threshold $\epsilon$, the prey becomes perceptually indistinguishable from the background. The predator’s perceptual system has not abandoned physics; it has been deceived at the level of compressed invariants.

This dynamic reveals a crucial insight. Seeming can simulate being if it replicates the invariants recognized by a compressed perceptual system. Structural grounding is not required if the signal triggers invariant detection heuristics. Camouflage does not alter the substrate of the organism; it alters the feature space in which recognition occurs.

The spectacle may be interpreted as a socio-cognitive analogue of camouflage. It produces signals optimized to match the invariant templates embedded in collective perception. Signals associated with authenticity, excellence, partnership, or community are staged with high salience. Yet the underlying structural substrate may diverge from these signals. Recognition is triggered; structural reciprocity is absent.

In this framework, being corresponds to transformation-invariant structure across physical and social dimensions. Seeming corresponds to the simulation of those invariants at the level of compressed representation. The motto \textit{Esse quam videri} thus acquires a physiological interpretation. To be rather than to seem is to maintain invariant coherence under transformation rather than merely to trigger recognition patterns.

The spectacle is not flawed because it produces images. All perception involves mediation. It is flawed when it systematically amplifies invariant-simulating signals while severing correspondence with substrate constraints. In doing so, it exploits the very mechanisms by which organisms learn physics. It converts the efficiency of sparse inference into an opening for large-scale distortion.

To move from analogy to formal critique, we must specify how divergence between signal and substrate can be modeled. The next section introduces a functional measure of spectacular distortion and examines how energetic flows distinguish mutual integration from parasitic extraction.

\section{Formalizing the Spectacle: Invariant Structure and Distortion}

Having grounded the distinction between being and seeming in perceptual invariance and sparse coding, we can now formalize the spectacle as a structural divergence between signal and substrate. The goal is not to reduce social theory to mathematics, but to articulate a constraint that renders the critique of appearing measurable in principle.

Let $B$ denote a structural substrate. In economic terms, $B$ may represent labor flows, energetic resources, institutional constraints, or material conditions. Let $S$ denote a social signal. $S$ may take the form of branding, reputation, narrative framing, symbolic capital, or public image. In a structurally congruent system, $S$ is a projection of $B$ that preserves invariant features relevant to collective inference.

We may formalize this relationship by introducing a mapping $\sigma : B \to S$ such that $S = \sigma(B)$. In a healthy system, $\sigma$ preserves invariant structure under relevant transformation groups. That is, if $g$ belongs to a transformation group $G$ acting on $B$, then
\[
\sigma(g \cdot B) \approx g' \cdot \sigma(B),
\]
for some corresponding transformation $g'$ in the signal space. The signal tracks the structural dynamics of its substrate.

To model distortion, let $\phi$ denote the compressed perceptual encoding applied by collective cognition to both $B$ and $S$. We define a divergence functional
\[
D(S, B) = \|\phi(S) - \phi(\sigma(B))\|.
\]
When $D(S,B)$ approaches zero, signal and substrate are congruent at the level of recognized invariants. Appearance tracks being. When $D(S,B)$ exceeds a critical threshold, signals simulate invariant structure without preserving substrate coherence. Seeming displaces being.

This formulation reframes Debord’s claim that the spectacle replaces lived reality with representation. The replacement is not total in ontological terms; the substrate continues to exist. Rather, the replacement occurs in the perceptual field. Compressed representations converge on $S$ while diverging from $B$. Collective inference stabilizes around image rather than structure.

The energetic dimension clarifies the difference between mutual integration and parasitic extraction. Consider a host system with energetic capacity $E_h$ and a platform or institutional structure with extraction $E_p$. In endosymbiosis, integration increases total system capacity such that
\[
\Delta E_{\text{total}} > 0 \quad \text{and} \quad \Delta E_h > 0.
\]
Structural congruence is maintained; signals of partnership correspond to mutual gain. In parasitism,
\[
\Delta E_h < 0,
\]
even if surface signals encode narratives of empowerment, collaboration, or opportunity. The signal space $S$ preserves invariants associated with symbiosis, while the substrate $B$ reflects asymmetric extraction. The divergence functional $D(S,B)$ increases.

Competition provides another illustrative case. If competitive signaling corresponds to genuine gradients of skill and contribution, then invariant features in $S$ track structural performance in $B$. However, if competitive framing masks precarity, instability, or artificial scarcity engineered for attention extraction, then the signal simulates merit without structural reciprocity. The spectacle thus rewards invariant-simulating surfaces.

Predictive processing deepens this account. Cognitive systems maintain priors about invariant structure in social environments. High-salience cues that match these priors reduce local prediction error. The spectacle supplies precisely such cues. It presents compressed patterns that align with templates for authenticity, excellence, or belonging. At the same time, it suppresses error signals that would reveal substrate divergence. The system achieves perceptual stability while accumulating structural incoherence.

The spectacle, therefore, is not a mere accumulation of images. It is a regime in which $D(S,B)$ is systematically amplified while perceptual compression prevents its detection. It exploits the fact that organisms must rely on sparse feature sets to navigate complex environments. What Debord described in moral and historical terms can be reframed as a regime of high signal–substrate divergence.

The injunction \textit{Esse quam videri} becomes, under this model, a constraint condition:
\[
D(S,B) \to 0.
\]
To be rather than to seem is to maintain minimal divergence between invariant substrate and perceptual signal. It is not a rejection of appearance. Appearance is unavoidable. Rather, it is a demand that appearance preserve transformation-invariant structure across energetic and institutional dynamics.

The spectacle becomes pathological when it normalizes high divergence as standard operating condition. Under such circumstances, seeming is optimized for recognition, while being erodes under extraction. The critique shifts from authenticity nostalgia to structural congruence. The question is not whether representation exists, but whether representation preserves invariant coherence.

\section{The Wizard and Sprezzatura: Illusion, Effortlessness, and Concealed Generation}

Two cultural motifs illuminate the dynamics of seeming without being: the Wizard of Oz and the Renaissance ideal of \textit{sprezzatura}. Both revolve around the management of appearance, yet they differ in how they relate surface display to generating structure.

In \textit{The Wonderful Wizard of Oz}, the central revelation is not that the Wizard possesses hidden supernatural power, but that he possesses none. The grand projections of authority—flames, booming voices, mechanized spectacle—are stage effects. Behind the curtain stands an ordinary man operating a machine. The spectacle is sustained by concealment of its generating apparatus. The visible implementation is optimized to trigger awe, fear, and submission. The generating function is deliberately hidden.

The Wizard is therefore a paradigmatic case of high divergence between signal and substrate. The signal space $S$ encodes omnipotence. The substrate $B$ consists of mechanical props and rhetorical manipulation. The divergence functional $D(S,B)$ is large, yet recognition persists because the audience’s perceptual compression relies on salient cues rather than causal inspection. The spectacle works until the curtain is drawn.

The Wizard illustrates that appearing can be engineered through control of perceptual channels. The lesson is not merely moral; it is structural. When the generating mechanism is concealed and insulated from feedback, the signal can be amplified beyond substrate capacity. Authority becomes theatrical rather than invariant.

By contrast, the Renaissance concept of \textit{sprezzatura}, articulated by Baldassare Castiglione in \textit{Il Cortegiano}, presents a subtler phenomenon. Sprezzatura denotes cultivated effortlessness. It is the art of performing skill in such a way that labor disappears from view. The generating structure—discipline, rehearsal, technical mastery—remains real. What is concealed is not emptiness, but effort.

Sprezzatura thus occupies an intermediate position between being and seeming. The substrate $B$ is robust. The performer genuinely possesses the invariant structure required for excellence. However, the signal $S$ suppresses visible strain. The performance appears natural, spontaneous, unforced. The generating function remains intact, but its process is aesthetically obscured.

Formally, in the Wizard case, $B$ lacks the structure implied by $S$. In sprezzatura, $B$ exceeds what $S$ explicitly displays. The divergence functional behaves differently. In the Wizard scenario, $D(S,B)$ is high because invariant structure is absent. In sprezzatura, $D(S,B)$ may remain low with respect to outcome invariants, yet high with respect to visible process. The concealment concerns implementation cost, not structural capacity.

This distinction clarifies an important nuance. Not all concealment is parasitic. The spectacle becomes pathological when signals simulate invariant structure without substrate grounding. Sprezzatura, by contrast, conceals labor while preserving invariant competence. It is a management of surface without falsifying the generating principle.

Nevertheless, sprezzatura risks drift toward spectacle if optimization of surface becomes detached from substrate cultivation. When the aesthetic of effortlessness is pursued independently of disciplined formation, the practice collapses into Wizardry. The performer maintains the signal of mastery without sustaining the generating function. Recognition persists temporarily, but structural erosion follows.

Both motifs therefore reinforce the methodological reading of \textit{Esse quam videri}. The Wizard warns against mistaking amplified signal for generating power. Sprezzatura demonstrates that surface modulation need not violate structural congruence, provided that invariant substrate remains intact. The critical distinction lies not in whether effort is visible, but in whether invariant generating structure exists.

In a society saturated by spectacle and profilicity, the temptation is to optimize the curtain rather than the mechanism. The discipline required is to inquire into the generating function. Does the visible implementation reflect stable invariants under transformation? Or is it a stage effect sustained by compression limits and attention control?

The Wizard dramatizes divergence. Sprezzatura dramatizes disciplined concealment. Together, they show that seeming is not inherently false, but that it becomes deceptive when the generating structure is hollow. The task remains the same: to preserve congruence between signal and substrate, even when the curtain is drawn.

The next section turns to a contemporary framework that rejects Debord’s narrative of lost unity while retaining the insight that identity and value are mediated through images. The concept of profilicity reframes the problem not as fall from authenticity, but as transformation in identity technology.

\section{Profilicity: Identity After Authenticity}

The critique of the spectacle culminates, in Debord’s formulation, in a narrative of lost authenticity. He suggests that a prior unity of life has been fractured by representation. What was once directly lived has become mediated display. The revolutionary impulse in his work is animated by the desire to restore immediacy, to abolish separation, and to reclaim an authentic mode of being.

A contemporary alternative challenges this nostalgia. The theory of profilicity argues that identity has never been grounded in pure, unmediated presence. Historical existence has always been structured by incongruence, role differentiation, and social evaluation. The difference in the present is not the emergence of mediation, but the dominance of a particular identity technology: the profile.

Profilicity names a regime in which individuals are constituted through publicly visible representations validated by a generalized peer. The profile is not simply a curated persona on a digital platform. It is a structural form in which identity becomes legible through metrics, feedback loops, and comparative visibility. Recognition is distributed through systems of ratings, likes, citations, rankings, and endorsements. Value emerges relationally within a network of mutual observation.

From the standpoint of profilicity, the distinction between being and seeming is not eliminated, but reconfigured. The profile is a constructed image, yet it is not necessarily an illusion masking a hidden authentic core. Rather, it is a functional interface through which social participation occurs. Identity is technologically scaffolded. There is no original unity to which one might return; there are only successive regimes of mediation.

This reframing modifies the interpretation of divergence. Under Debord’s model, high $D(S,B)$ signals pathological displacement of reality by image. Under profilicity, the signal $S$ is itself part of the substrate. Profiles generate real consequences. They shape opportunities, access, and social mobility. The mapping $\sigma : B \to S$ becomes bidirectional. Signals feed back into substrate conditions, altering energetic flows and institutional structures.

The risk, however, remains. Profilicity does not abolish divergence. It redistributes it. When identity becomes dependent on validation from a generalized peer, the incentive to optimize $S$ independently of $B$ intensifies. Signals are refined to trigger recognition heuristics. The compression constraints described earlier still apply. Sparse coding at the collective level privileges salient features over structural depth. The vulnerability to camouflage persists.

The crucial difference is conceptual rather than structural. Profilicity rejects authenticity as a lost historical state. It treats identity as technologically mediated without presupposing a pre-spectacular essence. In this view, critique must operate from within the regime of mediation rather than from an imagined exterior. One cannot step outside profiles; one can only modulate the degree of divergence they introduce.

The demand for structural congruence thus shifts from restoration to calibration. The question becomes whether profiles preserve invariant structure across transformation or whether they amplify signal–substrate divergence for competitive advantage. Being is not redefined as pre-representational presence, but as coherence under iterative feedback.

In this light, the tension introduced by \textit{Esse quam videri} remains operative but loses its nostalgic dimension. To be rather than to seem does not imply withdrawal from visibility. It implies maintaining low divergence between encoded identity and energetic reality within a mediated system. The problem is not that we inhabit a world of profiles. The problem arises when profiles simulate invariants that the substrate cannot sustain.

The spectacle and profilicity therefore share a semiotic foundation while diverging in historical interpretation. Both acknowledge that value and meaning circulate through images. Debord interprets this as a fall from unity; profilicity interprets it as an identity regime. The physiological and mathematical framework introduced here does not adjudicate between nostalgia and acceptance. It introduces a structural metric: divergence.

Under that metric, critique becomes continuous rather than revolutionary. The aim is not to abolish mediation, but to minimize distortion. The self is inevitably profiled. The task is to ensure that the profile does not become a camouflage layer severed from substrate coherence.


\section{Abstraction and the Disappearance of Implementation}

In any sufficiently developed system, abstraction operates by reducing irrelevant degrees of freedom. A model that preserves every micro-detail is unusable. What matters is the retention of structure sufficient to generate stable behavior across transformation. The distinction between being and seeming can therefore be reinterpreted as a question about abstraction fidelity.

When abstraction functions properly, it preserves the generating constraints of a system while suppressing incidental variation. A differential equation abstracts from specific trajectories while retaining the law that produces them. A geometric invariant abstracts from coordinate choice while preserving relational structure. The abstract representation is not a distortion; it is a compression that maintains causal coherence.

Pathology arises when abstraction ceases to preserve the generating constraint and instead optimizes for perceptual salience. In such cases, the abstract layer floats free of the substrate it once summarized. The visible interface becomes self-referential. The implementation disappears not because it is unnecessary, but because it has been displaced by a representation optimized for recognition rather than generation.

Under this condition, systems become evaluated according to their surface consistency rather than their structural resilience. Stability under scrutiny is replaced by stability under attention. The question shifts from whether a rule generates coherent outcomes to whether its projection triggers the correct response. Abstraction becomes camouflage.

\section{Performance and the Illusion of Causality}

Consider systems in which outcomes appear effortless, instantaneous, or inevitable. The disappearance of visible process can be interpreted in two ways. In one case, the underlying generative mechanism is intact and robust; the suppression of visible effort reflects mastery. In another case, the appearance of inevitability conceals the absence of generative depth. The surface is maintained through control of perceptual channels rather than through structural capacity.

The illusion of causality arises when visible signals are mistaken for the forces that generate them. Metrics are taken as causes rather than as indicators. Reputation becomes conflated with competence. Visibility becomes conflated with value. In each instance, the system confuses projection with principle.

This confusion can be formalized as a reversal of generative hierarchy. Instead of deriving instances from a stable rule, the system derives evaluation from surface manifestation. The particular replaces the general. Implementation displaces generation. What appears to function becomes authoritative irrespective of whether it can reproduce itself under perturbation.

\section{Structural Testing Under Transformation}

A reliable indicator of being, as distinct from seeming, is persistence under transformation. If a structure remains coherent when parameters shift, illumination changes, or context varies, then it exhibits invariant depth. If it collapses when attention wanes or when external conditions fluctuate, then it was sustained by surface optimization.

Testing under transformation provides a criterion independent of aesthetic impression. A theory remains valid if it generates consistent predictions across coordinate frames. An institution remains stable if it withstands resource fluctuation without rhetorical inflation. An identity remains coherent if it preserves relational commitments across audiences and temporal phases.

Seeming thrives in stable perceptual environments. Being endures perturbation. The difference lies not in intensity of display but in resistance to distortion. Systems optimized for surface recognition tend to degrade when the feedback loop shifts. Systems grounded in generating constraint maintain coherence even when visibility declines.

\section{The Discipline of Returning to Principle}

To interpret being as generating function is to orient analysis toward first principles rather than implementations. This does not entail rejection of implementation. Specific forms are necessary; interfaces are unavoidable. The question concerns priority. Does the system cultivate the rule from which instances arise, or does it curate instances while neglecting the rule?

In cognitive terms, this distinction parallels the difference between memorization and structural understanding. Memorization succeeds within a narrow distribution of cases. Structural understanding generalizes across perturbations. The former produces impressive surface fluency; the latter produces invariant competence.

The tension between being and seeming therefore becomes a question of generative depth. Systems that privilege projection over principle accumulate fragility. They must continually reinforce surface coherence because they lack invariant substrate. Systems that preserve generating constraints can tolerate variation without collapse.

The enduring relevance of \textit{Esse quam videri} lies in this structural insight. To be rather than to seem is to anchor implementation in generative law. It is to ensure that abstraction reduces without falsifying, that performance conceals labor without erasing structure, and that visibility remains a projection of invariant coherence rather than its substitute.

The final section synthesizes these threads and returns to the original injunction. It reframes the moral imperative as a structural constraint and situates critique within the condition it seeks to analyze.

\section{Conclusion: The Demand for Structural Congruence}

The phrase \textit{Esse quam videri} began as a school motto. It now stands as a structural constraint. The tension between being and seeming, first encountered as a moral exhortation, unfolds into a physiological, mathematical, and socio-economic problem. Guy Debord diagnosed a society in which appearing displaces being, and social relations are mediated by images. His critique articulated a condition of alienation in which representation becomes autonomous and authenticity is declared lost.

This essay has preserved the structural insight of that diagnosis while suspending its nostalgia. Rather than presupposing an original unity of life, it has grounded the distinction between being and seeming in perceptual invariance. Organisms learn physics by detecting transformation-invariant structure. Sparse coding enables efficiency but creates vulnerability. Camouflage exploits compression by simulating invariant features without preserving substrate reality. The spectacle can be understood as a large-scale analogue of such simulation.

By introducing the divergence functional
\[
D(S,B) = \|\phi(S) - \phi(\sigma(B))\|,
\]
the critique shifts from moral lament to formal constraint. Being corresponds to invariant coherence across transformations. Seeming corresponds to invariant simulation at the level of compressed representation. When $D(S,B)$ approaches zero, appearance tracks structure. When $D(S,B)$ increases, recognition is triggered without reciprocity. The spectacle becomes pathological not because it produces images, but because it amplifies divergence between signal and substrate.

The comparison with profilicity further clarifies the stakes. If identity is technologically mediated and there is no pristine past to restore, then critique must operate from within mediation. Profiles are not inherently illusions; they are interfaces. The risk lies in optimizing signal independent of structural grounding. The task is not to abolish appearance, but to calibrate it.

Structural congruence therefore replaces authenticity nostalgia. The injunction to be rather than to seem becomes a demand for low divergence under transformation. It requires that energetic flows, institutional structures, and identity signals remain reciprocally aligned. It requires that the invariants detected by collective perception correspond to constraints sustained by the substrate.

Such a demand cannot be satisfied through withdrawal from visibility. Nor can it be achieved by revolutionary rupture alone. It requires continuous self-critical awareness. Every critique risks becoming another signal within the perceptual field it interrogates. Every posture against spectacle may itself become spectacular. The divergence functional applies not only to institutions but to theorists.

The authentic journey, if the term is retained, does not consist in returning to pre-representational immediacy. It consists in preserving coherence between signal and structure across time. It consists in ensuring that appearing remains a projection of being rather than its camouflage. In a world where compression is inevitable and mediation is structural, the preservation of invariant congruence becomes the central ethical and epistemic task.


\newpage
\section*{Appendices}

\appendix

\section{Appendix A: Transformation Groups and Invariant Structure}

Let $X$ be a smooth manifold representing environmental or structural state space. Let $G$ be a Lie group acting smoothly on $X$:
\[
\Phi : G \times X \to X, \quad (g,x) \mapsto g \cdot x.
\]

A function $f : X \to \mathbb{R}^n$ is $G$–invariant if
\[
f(g \cdot x) = f(x), \quad \forall g \in G, \; \forall x \in X.
\]

Define the orbit of $x \in X$:
\[
\mathcal{O}_x = \{ g \cdot x \mid g \in G \}.
\]

An invariant $f$ is constant on orbits:
\[
f|_{\mathcal{O}_x} = \text{const}.
\]

Let $T_xX$ denote the tangent space at $x$. The infinitesimal action of the Lie algebra $\mathfrak{g}$ generates vector fields
\[
\xi_X(x) = \frac{d}{dt}\bigg|_{t=0} \exp(t\xi) \cdot x.
\]

Invariance condition (infinitesimal form):
\[
\mathcal{L}_{\xi_X} f = 0, \quad \forall \xi \in \mathfrak{g}.
\]

Structural being is identified with equivalence classes under $G$:
\[
[x] = X / G.
\]

Projection to quotient:
\[
\pi : X \to X/G.
\]

Invariant structure corresponds to functions descending to the quotient:
\[
f = \tilde{f} \circ \pi.
\]

\section{Appendix B: Sparse Encoding and Compressed Representation}

Let $X \subset \mathbb{R}^m$ denote high-dimensional state space.  
Let $Z \subset \mathbb{R}^k$ with $k \ll m$ denote compressed representation space.

Define encoding map
\[
\phi : X \to Z.
\]

Assume dictionary matrix $B = [b_1 \dots b_n] \in \mathbb{R}^{m \times n}$.

Sparse coding solves
\[
\min_{a \in \mathbb{R}^n} \|x - Ba\|_2^2 + \lambda \|a\|_1.
\]

Compressed representation:
\[
\phi(x) = a^*(x),
\]
where $a^*(x)$ is minimizer of the above functional.

Discrimination condition between $x,y \in X$:
\[
\|\phi(x) - \phi(y)\|_2 > \epsilon.
\]

Camouflage condition:
\[
\|\phi(x) - \phi(y)\|_2 \le \epsilon.
\]

If $G$ acts on $X$, define equivariance defect:
\[
E(g,x) = \|\phi(g \cdot x) - \rho(g)\phi(x)\|_2,
\]
where $\rho$ is representation of $G$ on $Z$.

Perfect equivariance:
\[
E(g,x)=0.
\]

Compression vulnerability:
\[
\exists x \ne y \text{ such that } \phi(x)=\phi(y).
\]

\section{Appendix C: Signal–Substrate Divergence Functional}

Let $B \subset \mathbb{R}^m$ denote substrate space.  
Let $S \subset \mathbb{R}^n$ denote signal space.

Let $\sigma : B \to S$ be projection map.

Let $\phi_S : S \to Z$, $\phi_B : B \to Z$ be compressed encodings into common feature space $Z$.

Define divergence functional
\[
D(S,B) = \|\phi_S(S) - \phi_B(B)\|_2.
\]

Expectation over distribution $\mu$ on $B$:
\[
\mathcal{D} = \mathbb{E}_{B \sim \mu} \left[ D(\sigma(B),B) \right].
\]

Congruence condition:
\[
\mathcal{D} \to 0.
\]

Distortion regime:
\[
\mathcal{D} > \delta,
\]
for threshold $\delta > 0$.

Let energetic functional $E : B \to \mathbb{R}_{\ge 0}$.

Signal consistency constraint:
\[
\nabla_S E(\sigma^{-1}(S)) \approx 0.
\]

Parasitic extraction condition:
\[
\frac{d}{dt}E(B(t)) < 0
\quad \text{while} \quad
\frac{d}{dt}\|\phi_S(S(t))\|_2 \ge 0.
\]

Stability under perturbation $\eta$:
\[
D(\sigma(B+\eta),B+\eta) - D(\sigma(B),B) = O(\|\eta\|).
\]

\section{Appendix D: Generating Operators and Implementation Maps}

Let $\Theta \subset \mathbb{R}^p$ denote parameter space.  
Let $G : \Theta \to X$ be generating operator.

Observable instance:
\[
x = G(\theta).
\]

Let $\Pi : X \to S$ be implementation map to signal space.

Composite mapping:
\[
\Theta \xrightarrow{G} X \xrightarrow{\Pi} S.
\]

Structural fidelity condition:
\[
\mathrm{rank}(DG_\theta) = \mathrm{rank}(D(\Pi \circ G)_\theta),
\]
where $D$ denotes differential.

Loss of generative depth:
\[
\mathrm{rank}(D(\Pi \circ G)_\theta) < \mathrm{rank}(DG_\theta).
\]

Overfitting to implementation:
\[
\exists \theta_1 \ne \theta_2 \text{ such that }
\Pi(G(\theta_1)) = \Pi(G(\theta_2)).
\]

Equivalence relation induced by implementation:
\[
\theta_1 \sim \theta_2
\iff
\Pi(G(\theta_1)) = \Pi(G(\theta_2)).
\]

Generative invariance under transformation group $H$ acting on $\Theta$:
\[
G(h \cdot \theta) = g_h \cdot G(\theta),
\]
for induced $g_h \in G$ acting on $X$.

Projection stability condition:
\[
\Pi(g_h \cdot x) = \rho(h)\Pi(x),
\]
for representation $\rho$ of $H$ on $S$.

Implementation distortion functional:
\[
\Delta(\theta)
=
\| DG_\theta - D(\Pi \circ G)_\theta \|.
\]

\section{Appendix E: Predictive Processing and Error Suppression}

Let $X$ denote latent substrate states.  
Let $S$ denote observable signals.  

Generative model:
\[
p(S \mid X), \quad p(X).
\]

Predictive distribution:
\[
p(S) = \int p(S \mid X)p(X)\, dX.
\]

Prediction error functional:
\[
\mathcal{E}(S,X) = \| S - \hat{S}(X) \|_2^2,
\]
where $\hat{S}(X) = \mathbb{E}[S \mid X]$.

Variational free energy:
\[
\mathcal{F}(q,X)
=
\mathbb{E}_{q(X)}[\mathcal{E}(S,X)]
+
\mathrm{KL}(q(X)\|p(X)).
\]

Perceptual update:
\[
\dot{q}(X) = - \nabla_q \mathcal{F}.
\]

Signal optimization under perceptual compression:
\[
S^* = \arg\min_S \mathbb{E}_{q(X)}[\mathcal{E}(S,X)].
\]

Error suppression condition:
\[
\nabla_X \mathcal{E}(S,X) \approx 0
\quad \text{while} \quad
\|X - X_{\text{true}}\| \text{ large}.
\]

Structural misalignment regime:
\[
\mathcal{E}(S,X_{\text{true}}) \text{ small}
\quad \text{and} \quad
D(S,B) > \delta.
\]

Attractor dynamics:
\[
\dot{X} = - \nabla_X \mathcal{F}.
\]

Metastable signal basin:
\[
\nabla_X \mathcal{F} = 0,
\quad
\nabla_X^2 \mathcal{F} > 0.
\]

Persistent divergence:
\[
\exists \; \epsilon > 0
\text{ such that }
D(S(t),B(t)) > \epsilon
\quad \forall t \in [t_0,t_1].
\]

\section{Appendix F: Category-Theoretic Interpretation of Being, Seeming, and Divergence}

Let $\mathbf{B}$ be a category of substrates. Objects $b \in \mathrm{Ob}(\mathbf{B})$ represent structured conditions such as material processes, energetic configurations, institutional states, or any other generative bases. Morphisms $u:b\to b'$ represent lawful transformations of substrate, such as resource reallocations, policy updates, perturbations, or time evolution. Let $\mathbf{S}$ be a category of signals. Objects $s \in \mathrm{Ob}(\mathbf{S})$ represent representations, profiles, reputational states, or symbolic interfaces; morphisms $\alpha:s\to s'$ represent transformations of signal, such as rebranding, narrative revision, metric updates, or channel modulation.

A central primitive is a functor
\[
\Sigma : \mathbf{B} \to \mathbf{S},
\]
interpreted as a projection from substrate to signal. For each substrate object $b$, $\Sigma(b)$ is its represented appearance; for each substrate morphism $u:b\to b'$, $\Sigma(u):\Sigma(b)\to\Sigma(b')$ is the induced change in appearance. Functoriality enforces compositional consistency:
\[
\Sigma(\mathrm{id}_b)=\mathrm{id}_{\Sigma(b)}, 
\qquad 
\Sigma(v\circ u)=\Sigma(v)\circ \Sigma(u).
\]
In this language, a regime of representation is a choice of $\Sigma$.

The notion of invariance previously expressed via group actions is recast as categorical descent. Suppose a subcategory $\mathbf{G}\subseteq \mathrm{End}(\mathbf{B})$ encodes ``gauge'' transformations, i.e.\ morphisms that should not alter identity at the level of being. One may form a localization or quotient category $\mathbf{B}[\mathbf{G}^{-1}]$ in which the morphisms in $\mathbf{G}$ are inverted, or alternatively a categorical orbit construction when $\mathbf{G}$ acts as a groupoid. In either case, the structural content of being is identified with equivalence classes of objects under morphisms regarded as observationally irrelevant. A signal functor $\Sigma$ is said to respect invariance if it factors through the quotient:
\[
\mathbf{B} \xrightarrow{\pi} \mathbf{B}[\mathbf{G}^{-1}] \xrightarrow{\widetilde{\Sigma}} \mathbf{S},
\qquad
\Sigma = \widetilde{\Sigma}\circ \pi.
\]
This factorization expresses that representation depends only on invariant structure.

Compression is expressed by the failure of faithfulness. A functor $\Sigma$ is faithful if distinct substrate morphisms remain distinct after projection. A functor is full if every signal morphism between represented objects is induced by a substrate morphism. Compression in the strong sense appears as non-faithfulness: there exist $u\neq v : b\to b'$ such that $\Sigma(u)=\Sigma(v)$. This identifies distinct substrate transformations as identical at the level of appearance, which is the categorical analogue of feature-collision in sparse coding. Likewise, non-injectivity on objects expresses camouflage at the object level: there exist $b\not\cong b'$ with $\Sigma(b)\cong \Sigma(b')$.

A categorical articulation of ``seeming without being'' is therefore: the functor $\Sigma$ collapses structurally distinct substrate states into isomorphic or identical signal states, while the ambient evaluative processes occur in $\mathbf{S}$ rather than in $\mathbf{B}$. In this regime, social selection is performed on $\Sigma(b)$ rather than on $b$.

To model evaluation and recognition, introduce a category $\mathbf{E}$ of evaluations, such as rankings, decisions, allocations, and affordances. An evaluation functor
\[
\mathcal{V} : \mathbf{S} \to \mathbf{E}
\]
maps signals to evaluative outcomes. The composite
\[
\mathcal{V}\circ \Sigma : \mathbf{B} \to \mathbf{E}
\]
defines how substrate states are acted upon by passing through representation. A spectacle-like pathology can then be read as domination of $\mathcal{V}\circ\Sigma$ over any direct substrate-sensitive evaluation, meaning the effective dynamics of the system are governed by the composite rather than by intrinsic substrate constraints.

The divergence functional $D(S,B)$ can be categorified by introducing an ``invariant extractor'' functor. Let $\Phi_B:\mathbf{B}\to\mathbf{Z}$ and $\Phi_S:\mathbf{S}\to\mathbf{Z}$ be functors into a comparison category $\mathbf{Z}$ whose objects encode invariant descriptors, sufficient statistics, or any compressed features used for recognition. The earlier real-valued divergence is recovered by choosing $\mathbf{Z}$ enriched over metric spaces, but the categorical statement does not require metrics. Define a comparison natural transformation candidate
\[
\eta : \Phi_B \Rightarrow \Phi_S\circ \Sigma.
\]
When $\eta$ exists as a natural isomorphism, the representation preserves invariant structure in the strongest sense:
\[
\Phi_B(b) \cong \Phi_S(\Sigma(b))
\quad \text{naturally in } b.
\]
When such an $\eta$ fails to exist, or exists only as a lax transformation, the mismatch between invariant extraction from substrate and invariant extraction from signal formalizes the ``gap'' between being and seeming. In enriched settings one may define, for each object $b$, a distortion object or number measuring the failure of $\eta_b$ to be an isomorphism.

The distinction between generating principle and implementation admits a universal characterization. Let $\Theta$ be a category of parameters or causes and let $G:\Theta\to\mathbf{B}$ be a functor assigning to parameters their generated substrate states. The signal produced by implementation is $\Sigma\circ G:\Theta\to\mathbf{S}$. If one is given partial data on how appearances should behave, one can recover a best approximation to $G$ by Kan extension. Concretely, if $J:\Theta_0\hookrightarrow\Theta$ is an inclusion of observed parameter contexts and $G_0:\Theta_0\to\mathbf{B}$ is known, then the left Kan extension
\[
\mathrm{Lan}_J(G_0):\Theta\to\mathbf{B}
\]
is the universal way to extend the generating principle from observed contexts to the whole parameter space. A system optimized for implementation without principle corresponds to replacing $G$ by some $G'$ whose composite $\Sigma\circ G'$ matches observed appearances but does not satisfy the universal property of minimal extension in $\mathbf{B}$. In other words, it interpolates in $\mathbf{S}$ rather than generating in $\mathbf{B}$.

Energetic flow and endosymbiosis versus parasitism can be expressed using monoidal structure. Suppose $\mathbf{B}$ is symmetric monoidal with tensor product $\otimes$ representing composition of subsystems and unit object $\mathbb{1}$ representing the null system. An ``integration'' morphism $m:b\otimes c\to d$ represents merging two substrates into a composite. Introduce an energy valuation as a monoidal functor
\[
\mathcal{E}:\mathbf{B}\to (\mathbb{R}_{\ge 0}, +, 0)
\]
into the additive monoid viewed as a one-object strict monoidal category, or into an ordered monoidal poset. Endosymbiosis is then encoded by morphisms for which valuation is superadditive in the appropriate sense:
\[
\mathcal{E}(d) \ge \mathcal{E}(b)+\mathcal{E}(c),
\]
whereas parasitism corresponds to morphisms and subsequent dynamics in which one component’s valuation declines while the composite signal in $\mathbf{S}$ encodes an increase. The key point is that the monoidal structure makes ``mutualistic composition'' a structural statement rather than a narrative one.

A further refinement arises from factorization systems. Suppose $\mathbf{B}$ admits an (epi, mono)-factorization of morphisms, or more generally an orthogonal factorization system $(\mathcal{E},\mathcal{M})$. One can interpret $\mathcal{E}$ as ``generative'' or ``process'' morphisms that genuinely change substrate, and $\mathcal{M}$ as ``embedding'' morphisms that reorganize presentation without altering core structure. Seeming-dominant regimes correspond to dynamics in $\mathbf{S}$ that are dominated by images of $\mathcal{M}$-like morphisms, while being-dominant regimes correspond to dynamics that track images of $\mathcal{E}$-like morphisms. This expresses the intuition that some transformations are merely representational, while others are structurally causal.

Finally, the Wizard-versus-sprezzatura contrast admits a categorical formulation. The Wizard corresponds to a functor $\Sigma$ that is not conservative: it maps a non-isomorphism in $\mathbf{B}$ (lack of power) to an isomorphism-like pattern in $\mathbf{S}$ (appearance of power), thereby failing to reflect isomorphisms. Sprezzatura corresponds to a conservative $\Sigma$ with additional structure that hides computational cost. One may model cost as a profunctor or as a functor into a cost-enriched category, and sprezzatura as the existence of a natural transformation that forgets cost while preserving correctness. Wizardry is the failure of correctness itself.

In summary, the categorical perspective replaces scalar measures with structural properties. Being is modeled by quotient-invariant content of $\mathbf{B}$; seeming is modeled by the image in $\mathbf{S}$; compression is non-faithfulness; camouflage is object- and morphism-level collapse; congruence is the existence of a natural isomorphism between invariant extractors; distortion is the obstruction to such naturality; principle is universal generation via Kan extension; and parasitism versus symbiosis is expressed through monoidal valuation and functorial coherence. This formalism is deliberately general: it does not prescribe a single implementation, but isolates the generating constraints under which representations either preserve or sever structural invariants.

\newpage
\begin{thebibliography}{99}

\bibitem{Barthes1957}
Barthes, Roland.
\textit{Mythologies}.
Paris: Éditions du Seuil, 1957.

\bibitem{Baudrillard1981}
Baudrillard, Jean.
\textit{Simulacres et Simulation}.
Paris: Éditions Galilée, 1981.

\bibitem{Baum1900}
Baum, L. Frank.
\textit{The Wonderful Wizard of Oz}.
Chicago: George M. Hill Company, 1900.

\bibitem{Bateson1972}
Bateson, Gregory.
\textit{Steps to an Ecology of Mind}.
Chicago: University of Chicago Press, 1972.

\bibitem{Benjamin1936}
Benjamin, Walter.
“The Work of Art in the Age of Mechanical Reproduction.”
1936.

\bibitem{Bourdieu1984}
Bourdieu, Pierre.
\textit{Distinction: A Social Critique of the Judgement of Taste}.
Cambridge, MA: Harvard University Press, 1984.

\bibitem{Castiglione1528}
Castiglione, Baldassare.
\textit{Il Libro del Cortegiano}.
Venice: Aldus Manutius, 1528.

\bibitem{CiceroLaelius}
Cicero, Marcus Tullius.
\textit{Laelius de Amicitia}.
44 BCE.

\bibitem{Debord1967}
Debord, Guy.
\textit{La Société du Spectacle}.
Paris: Buchet-Chastel, 1967.

\bibitem{Debord1973}
Debord, Guy.
\textit{La Société du Spectacle} (Film).
Paris, 1973.

\bibitem{Derrida1967}
Derrida, Jacques.
\textit{De la grammatologie}.
Paris: Éditions de Minuit, 1967.

\bibitem{Eco1976}
Eco, Umberto.
\textit{A Theory of Semiotics}.
Bloomington: Indiana University Press, 1976.

\bibitem{Feuerbach1841}
Feuerbach, Ludwig.
\textit{Das Wesen des Christentums}.
Leipzig: Wigand, 1841.

\bibitem{Foucault1969}
Foucault, Michel.
\textit{L’archéologie du savoir}.
Paris: Gallimard, 1969.

\bibitem{Foucault1975}
Foucault, Michel.
\textit{Surveiller et punir}.
Paris: Gallimard, 1975.

\bibitem{Friston2010}
Friston, Karl.
“The free-energy principle: A unified brain theory?”
\textit{Nature Reviews Neuroscience}, 11(2):127–138, 2010.

\bibitem{Gentzen1935}
Gentzen, Gerhard.
“Untersuchungen über das logische Schließen.”
\textit{Mathematische Zeitschrift} 39 (1935): 176–210, 405–431.

\bibitem{Gibson1979}
Gibson, James J.
\textit{The Ecological Approach to Visual Perception}.
Boston: Houghton Mifflin, 1979.

\bibitem{Grothendieck1960}
Grothendieck, Alexander.
Éléments de géométrie algébrique.
\textit{Publications Mathématiques de l'IHÉS}, 1960–1967.

\bibitem{Heidegger1927}
Heidegger, Martin.
\textit{Sein und Zeit}.
Tübingen: Niemeyer, 1927.

\bibitem{Heidegger1954}
Heidegger, Martin.
\textit{Die Frage nach der Technik}.
Pfullingen: Neske, 1954.

\bibitem{Hinton2006}
Hinton, Geoffrey E., Simon Osindero, and Yee-Whye Teh.
“A fast learning algorithm for deep belief nets.”
\textit{Neural Computation}, 18(7):1527–1554, 2006.

\bibitem{Husserl1913}
Husserl, Edmund.
\textit{Ideen zu einer reinen Phänomenologie und phänomenologischen Philosophie}.
1913.

\bibitem{KashiwaraSchapira1990}
Kashiwara, Masaki, and Pierre Schapira.
\textit{Sheaves on Manifolds}.
Berlin: Springer, 1990.

\bibitem{Lacan1966}
Lacan, Jacques.
\textit{Écrits}.
Paris: Éditions du Seuil, 1966.

\bibitem{Luhmann1984}
Luhmann, Niklas.
\textit{Soziale Systeme}.
Frankfurt am Main: Suhrkamp, 1984.

\bibitem{MacLane1998}
Mac Lane, Saunders.
\textit{Categories for the Working Mathematician}.
2nd ed., New York: Springer, 1998.

\bibitem{Marx1844}
Marx, Karl.
\textit{Economic and Philosophic Manuscripts of 1844}.
1844.

\bibitem{Marx1867}
Marx, Karl.
\textit{Das Kapital}, Vol. I.
Hamburg: Otto Meissner, 1867.

\bibitem{MerleauPonty1945}
Merleau-Ponty, Maurice.
\textit{Phénoménologie de la perception}.
Paris: Gallimard, 1945.

\bibitem{Moeller2021}
Moeller, Hans-Georg, and Paul J. D’Ambrosio.
\textit{You and Your Profile: Identity After Authenticity}.
New York: Columbia University Press, 2021.

\bibitem{MoellerVideo}
Moeller, Hans-Georg.
“Guy Debord and the Society of the Spectacle.”
Video lecture transcript.

\bibitem{Reynolds1983}
Reynolds, John C.
“Types, Abstraction, and Parametric Polymorphism.”
\textit{Information Processing} 83 (1983): 513–523.

\bibitem{Samuel}
The Holy Bible.
1 Samuel 16:7.
Various editions.

\bibitem{Shannon1948}
Shannon, Claude E.
“A mathematical theory of communication.”
\textit{Bell System Technical Journal}, 27(3):379–423, 1948.

\bibitem{Spivak2014}
Spivak, David I.
\textit{Category Theory for the Sciences}.
Cambridge, MA: MIT Press, 2014.

\bibitem{Todorov1984}
Todorov, Tzvetan.
\textit{Theories of the Symbol}.
Ithaca: Cornell University Press, 1984.

\bibitem{Tomasello1999}
Tomasello, Michael.
\textit{The Cultural Origins of Human Cognition}.
Cambridge, MA: Harvard University Press, 1999.

\bibitem{Turkle2011}
Turkle, Sherry.
\textit{Alone Together: Why We Expect More from Technology and Less from Each Other}.
New York: Basic Books, 2011.

\bibitem{Wadler1989}
Wadler, Philip.
“Theorems for Free!”
In \textit{Proceedings of the Fourth International Conference on Functional Programming Languages and Computer Architecture}, 1989.

\bibitem{Weick1995}
Weick, Karl E.
\textit{Sensemaking in Organizations}.
Thousand Oaks: Sage Publications, 1995.

\end{thebibliography}

\end{document}